
\begin{proposition}[A single common cubic zero in the rational double-cover case]\label{cdp:st:C.2.RAT.8}
Retain the torus-fibre, rational-normalization and specified
ruled-boundary morphism of \ref{cdp:st:C.2.RAT.1}. Suppose $C$ is not
contracted, its image has generic ordinary-triple transversal,
and $C\to H^\nu$ has degree two. If the common zero divisor
of the cubic normal coefficients has degree one, then
$e(S)\ge0$.
\end{proposition}
\begin{proof}
As in \ref{cdp:st:C.2.RAT.7}, $H$ is smooth rational. On $W=\PP(N)$
write the complete root divisor as
\[
 B=R_1+R_2+F\sim3\xi,
 \qquad B_h=R_1+R_2\sim3\xi-F.
\tag{C.RD.17}\label{cdp:eq:C.RD.17}
\]
Here $F$ is the ruling fibre over the unique simple
common zero. The horizontal divisor $B_h$ is finite
over $H$, $R_1$ is a section, and the normalization of
$R_2$ is the rational degree-two cover $C$.
Since $K_W=-2\xi$, $\xi^2=2$, and $\xi F=1$,
\[
 p_a(B_h)=2,
 \qquad \sum_p\delta_p(B_h)=3,
 \qquad I(F,B_h)=3.
\tag{C.RD.18}\label{cdp:eq:C.RD.18}
\]
The full divisor $B$ has total delta six. We will
use its particular decomposition \eqref{cdp:eq:C.RD.17}, not a
statement about arbitrary delta-six plane germs.

\paragraph{The local cubic-suspension genus along $F$.}
At a point of $F$, put $i=I_p(F,B_h)$. If $i=1$,
the germ of $B$ is an ordinary node. If $i=2$,
the horizontal germ is either smooth with contact
two with $F$, two smooth branches transverse to
$F$, or a multiplicity-two cusp transverse to $F$.
The bound $\delta_p(B_h)\le3$ makes the resulting
plane types respectively
\[
 A_3,\qquad D_{2k+2}\ (1\le k\le3),\qquad
                 D_{2k+3}\ (1\le k\le3).
\tag{C.RD.19}\label{cdp:eq:C.RD.19}
\]
For two smooth branches $k$ is their contact;
for the cusp $k$ is its delta invariant.
The exact analytic identifications are the graph
straightening and square-completion comparisons
in \ref{cdp:st:C.2.RAT.6}. Their cubic suspensions have
geometric genus zero or one by (C.RD.10).

If $i=3$, the point contains $R_1$ and all the
intersection of $R_2$ with $F$. The section $R_1$
is transverse to $F$. A smooth branch of $R_2$
with intersection two is tangent to $F$ and
transverse to $R_1$, giving $D_6$. If $R_2$ has
two local branches, both are smooth and transverse
to $F$. The horizontal delta bound three forces
all three horizontal branches to be pairwise
transverse. The full germ is therefore four
distinct transverse smooth branches.
The last possibility is one singular branch
of $R_2$. Its multiplicity is two, since its
intersection with $F$ is two. Its intersection
with $R_1$ is at least two; hence \eqref{cdp:eq:C.RD.18}
forces it to be an ordinary cusp, transverse
to both smooth branches $F$ and $R_1$.
These alternatives exhaust $i=3$.

For four transverse branches, the degree-four
initial plane polynomial has four distinct
linear factors. Its cubic suspension is Newton
nondegenerate with weights $(1/4,1/4,1/3)$;
the only positive lattice point on or below
its compact Newton face is $(1,1,1)$. Its
geometric genus is therefore one by
\cite[Theorem 2.1.1 and Corollary 2.1.3]{MerleTeissier1977}.

We give the precise same calculation for the cusp
and two transverse branches, retaining higher
terms of the actual equation. Straighten $F,R_1$
to coordinate axes and make a linear change so
the cusp tangent is $u=0$. Up to units, the
equation is then
\[
 r(u,v)=v(u-v)\bigl(u^2+c v^3+
           u\,O((u,v)^2)+O((u,v)^4)\bigr),
                  \qquad c\ne0.
\tag{C.RD.20}\label{cdp:eq:C.RD.20}
\]
The nonzero radical cubic is the ordinary-cusp
condition. The compact Newton facets of
$z^3-r(u,v)$ have normalized inward vectors
\[
             (1/4,1/4,1/3),\qquad
             (3/10,1/5,1/3).
\tag{C.RD.21}\label{cdp:eq:C.RD.21}
\]
Their plane edges have vertex sets
$\{(3,1),(2,2)\}$ and
$\{(2,2),(0,5)\}$. The higher terms in
\eqref{cdp:eq:C.RD.20} are strictly above these faces.
The first edge polynomial is a nonzero multiple
of $u^2v(u-v)$; the second is a nonzero multiple
of $v^2(u^2+c v^3)$. Both have smooth zero sets
in $(\C^*)^2$. Each two-dimensional compact
face additionally contains the term $z^3$,
whose $z$-derivative is nonzero on $(\C^*)^3$.
This proves the required Newton nondegeneracy.
Every noncompact face inequality is strictly
satisfied by positive integral points; the only
additional one beyond coordinate positivity is
$v+z/3\ge1$.
For either vector in \eqref{cdp:eq:C.RD.21}, the point
$(1,1,1)$ has weight $5/6$, and increasing any
one coordinate makes its weight greater than
one. Thus it is the only positive lattice
point outside the interior of the Newton
polyhedron. The same adjunction theorem gives
$p_g(z^3-r)=1$.

This proves that every suspension singularity
over $F$ has geometric genus at most one.
A nonzero value requires $i\ge2$ and
$\delta_p(B_h)\ge1$: the cases with horizontal
delta zero were the ordinary node and $A_3$.
Since the total intersection with $F$ is three,
there is at most one such nonzero value over
$F$. Away from $F$, \ref{cdp:st:C.2.RAT.6} shows that a
nonzero geometric genus costs horizontal delta
at least three. It cannot coexist with a
nonzero value over $F$, by \eqref{cdp:eq:C.RD.18}.
Consequently the cyclic cover $X_0$ branched
over $B$ satisfies
\[
                        \sum_x p_g(X_0,x)\le1.
\tag{C.RD.22}\label{cdp:eq:C.RD.22}
\]

We also obtain the required link bound without
assuming a delta-six monodromy table. For a
normal surface singularity, let $E$ be the
reduced exceptional divisor of a good resolution.
The restriction maps from infinitesimal
neighbourhoods of $E$ onto $E$ are surjective
on $H^1$, since the kernels have one-dimensional
support and no $H^2$. Formal functions therefore
gives $h^1(E,\OO_E)\le p_g$.
For an SNC exceptional curve,
$b_1(E)=2\sum g_i+\beta\le
2(\sum g_i+\beta)=2h^1(E,\OO_E)$.
The invertible exceptional intersection matrix
identifies this $b_1(E)$ with the link's first
Betti number, as in the proof of \ref{cdp:st:C.2.RAT.6}.
Equation \eqref{cdp:eq:C.RD.22} yields
\[
              \lambda:=\sum_x b_1(\operatorname{link}(X_0,x))
                            \le2.
\tag{C.RD.23}\label{cdp:eq:C.RD.23}
\]

\paragraph{The compact pairing and the attachment count.}
The cyclic algebra and projective-bundle cohomology
are unchanged from (C.RD.13), so its projective
resolution $X$ has $b_1(X)\le2$.
The divisor $B$ is connected, since $F$ meets
both horizontal components. Its three classes
span a two-dimensional space: the section
$R_1$ and the fibre $F$ are independent.
Thus the boundary-relation number is exactly
$\varepsilon=1$, and $b_1(U)=b_1(X)+1\le3$.
Put $\nu=\sum_p(r_p(B)-1)$.
Normalization of its three rational components
and the resolved-boundary comparison give
\[
 e(U)=3(\nu-2),\qquad
 b_1(D)=\nu-2+\lambda,
 \qquad
 \operatorname{rank}I_U
       =2\nu-4+2b_1(X)-\lambda\ge2\nu-6.
\tag{C.RD.24}\label{cdp:eq:C.RD.24}
\]
The rank identity is the same Gysin, Hodge-index
and Euler calculation as (C.RD.15), with all
three boundary components and their one relation
included. Compact-core transport into the
nearby torus therefore forces $\nu\le3$.

Let $p\in H$ be the common zero value and
$a=\#(C\to H)^{-1}(p)\in\{1,2\}$.
At every point of $F\cap B_h$, the full divisor
has the branch of $F$ together with the
individual horizontal normalization branches.
It follows exactly that
\[
        \nu_F:=\sum_{x\in F}(r_x(B)-1)=1+a.
\tag{C.RD.25}\label{cdp:eq:C.RD.25}
\]
The contribution one is the section $R_1$;
the other $a$ branches are the normalization
points of $R_2$ above $p$. This identity retains
coinciding tangent roots and ramification.

The loss of distinct horizontal anchor values
at $p$ is at most $a-1=\nu_F-2$. Away from
$p$, the cubic is primitive. The argument at
the end of \ref{cdp:st:C.2.RAT.7} makes each duplicate
attachment there a multibranch point of $R_2$,
costing at least one in
$\nu_{\rm out}=\nu-\nu_F$. Hence the total
loss $\ell$ of distinct normalization-incidence
edges satisfies
\[
                    \ell\le\nu_F-2+\nu_{\rm out}
                              =\nu-2\le1.
\tag{C.RD.26}\label{cdp:eq:C.RD.26}
\]
Finally $r_{\rm dom}\le1$ and
$K\ge r-N_0-r_{\rm dom}$. Projecting these
$K$ surviving attachments to $H$ retains at
least $K-\ell$ distinct edges. Equation
(C.RA.24) and the Euler decomposition give
\[
                 e(S)\ge2-r_{\rm dom}-\ell\ge0.
\tag{C.RD.27}\label{cdp:eq:C.RD.27}
\]
The two Newton facets, their exact lower
lattice-point counts and the attachment
budgets are recorded in
\path{code/A_cubic_scalar_one.py} and its output.
\end{proof}


